\section{Skrótowa analiza wymagań}

Cel projektu: opracowanie metody (lub zestawu metod) pozwalającej odróżnić tekst napisany przez człowieka od tekstu wygenerowanego przez AI.

\begin{table}[H]
\centering
\begin{tabularx}{\textwidth}{|c|X|c|}
\hline
\textbf{ID} & \textbf{Wymaganie} & \textbf{Priorytet} \\
\hline
W1 & Klasyfikacja binarna: tekst ludzki / tekst AI & wysoki \\
\hline
W2 & Minimalna dokładność klasyfikacji $\geq 51\%$ na wybranym zbiorze testowym & wysoki \\
\hline
W3 & Obsługa tekstów o ustalonej minimalnej długości (do ustalenia eksperymentalnie) & średni \\
\hline
W4 & Porównanie co najmniej jednego podejścia statystycznego i jednego ML (o ile pozwolą dane) & średni \\
\hline
W5 & Wykorzystanie otwartego zbioru danych (lub własnego podzbioru) & wysoki \\
\hline
W6 & Udokumentowanie ograniczeń metody (krótkie teksty, domena, język) & średni \\
\hline
\end{tabularx}
\caption{Skrótowa lista wymagań}
\end{table}

Poprzeczka dokładności została świadomie ustawiona nisko ($51\%+$), ponieważ na etapie planowania nie jest pewne, czy stabilna detekcja jest możliwa przy dostępnych zasobach i danych - zwłaszcza dla języka polskiego.
